\documentclass[11pt,a4paper]{article}
\usepackage[utf8]{inputenc}
\usepackage[T1]{fontenc}
\usepackage[english]{babel}
\usepackage[margin=2.5cm]{geometry}
\usepackage{booktabs}
\usepackage{hyperref}
\usepackage{xcolor}
\usepackage{fancyhdr}
\usepackage{enumitem}
\usepackage{tabularx}
\usepackage{float}
\usepackage{listings}
\usepackage{tikz}
\usepackage{amsmath}
\usepackage{caption}

\definecolor{navy}{HTML}{1A3A8F}
\definecolor{dgray}{HTML}{6B7A99}
\definecolor{codebg}{HTML}{F4F6FA}
\definecolor{codeframe}{HTML}{C8D0E0}

\lstset{language=Python,basicstyle=\ttfamily\footnotesize,
  backgroundcolor=\color{codebg},frame=single,rulecolor=\color{codeframe},
  breaklines=true,keywordstyle=\bfseries\color{navy},
  commentstyle=\itshape\color{dgray},showstringspaces=false,
  xleftmargin=4pt,xrightmargin=4pt,aboveskip=6pt,belowskip=6pt,
  numbers=left,numberstyle=\tiny\color{dgray},numbersep=5pt}

\usepackage{titlesec}
\titleformat{\section}{\large\bfseries\color{navy}}{\thesection.}{0.5em}{}
\titleformat{\subsection}{\normalsize\bfseries\color{navy}}{\thesubsection}{0.5em}{}
\titleformat{\subsubsection}{\small\bfseries\itshape\color{dgray}}{\thesubsubsection}{0.5em}{}

\pagestyle{fancy}\fancyhf{}
\fancyhead[L]{\small\color{dgray}psgscoring v0.2.94 --- Handbook}
\fancyhead[R]{\small\color{dgray}psgscoring}
\fancyfoot[C]{\small\color{dgray}\thepage}
\renewcommand{\headrulewidth}{0.4pt}
\renewcommand{\footrulewidth}{0.2pt}

\title{\textbf{\LARGE psgscoring}\\[8pt]
\large A Technical Handbook for Automated\\
Polysomnography Scoring\\[8pt]
\normalsize Version 0.2.94 --- April 2026}
\author{psgscoring contributors\\
\url{https://github.com/bartromb/psgscoring}}
\date{}

\begin{document}
\maketitle
\thispagestyle{empty}

\begin{abstract}
\noindent
This handbook provides a complete technical introduction to
\texttt{psgscoring}, an open-source Python library (BSD-3) for
AASM~2.6-compliant automated polysomnography (PSG) scoring. It covers
clinical context, signal processing, event detection, apnea
classification, twelve bias corrections, hypoxic burden computation,
post-processing for central apnea refinement, and external validation.

\medskip\noindent
Prerequisites: Python~3.9+, NumPy, SciPy, basic DSP
(Fourier transform, filtering, envelopes).\\
Repository: \url{https://github.com/bartromb/psgscoring}\\
Installation: \texttt{pip install psgscoring}
\end{abstract}

\tableofcontents
\newpage

% ============================================================
\section{Clinical Context: Why This Matters}
% ============================================================

\subsection{Sleep Apnea}

Obstructive Sleep Apnea (OSA) affects approximately 936~million adults
worldwide. During sleep, the upper airway collapses partially or
completely, causing breathing interruptions that fragment sleep and
reduce blood oxygen levels. Left untreated, OSA increases the risk of
hypertension, stroke, heart failure, and traffic accidents.

Diagnosis relies on polysomnography (PSG): an overnight recording of
brain waves (EEG), eye movements (EOG), muscle tone (EMG), airflow,
respiratory effort, blood oxygen saturation (SpO$_2$), body position,
snoring, heart rhythm (ECG), and leg movements.

\subsection{The Scoring Problem}

A sleep technologist manually reviews the recording and counts:
\begin{itemize}[nosep]
  \item \textbf{Apneas}: $\geq$90\% flow reduction for $\geq$10~s
  \item \textbf{Hypopneas}: $\geq$30\% flow reduction for $\geq$10~s
    with $\geq$3\% oxygen desaturation or an EEG arousal
\end{itemize}
The Apnea-Hypopnea Index
$\text{AHI} = (\text{apneas} + \text{hypopneas})/\text{hours of sleep}$
is the primary diagnostic metric.
Manual scoring takes 1--2~hours per study with 10--20\% inter-scorer
variability. \texttt{psgscoring} automates this process.

\subsection{The AASM Standard}

The American Academy of Sleep Medicine (AASM) Manual for the Scoring of
Sleep and Associated Events (Version~2.6, 2020) is the rulebook
\cite{berry2020}. Key rules:
Rule~1A: hypopnea requires $\geq$3\% SpO$_2$ desaturation.
Rule~1B: hypopnea may alternatively be scored with an EEG arousal.
Apneas are scored on the thermistor; hypopneas on the nasal pressure
transducer.

\subsection{Beyond AHI: Why One Number Is Not Enough}

The AHI, though ubiquitous, has been critically appraised as a
reductionist metric. Pevernagie et~al.\ demonstrated that the
correlation between AHI and clinical endpoints---sleepiness,
hypertension, quality of life---is surprisingly weak, and that reliance
on the AHI leads to overestimation of OSA's clinical relevance
\cite{pevernagie2020}. Two patients with identical AHI values can
experience fundamentally different physiological consequences
\cite{pevernagie2025}.

\texttt{psgscoring} addresses this by computing three complementary
dimensions: (1)~AHI confidence interval (event frequency with
uncertainty), (2)~hypoxic burden (oxygen deficit per hour), and
(3)~arousal index (sleep fragmentation).

% ============================================================
\section{Architecture Overview}
% ============================================================

\texttt{psgscoring} is organised as 13 independently importable
submodules with a strict one-directional dependency graph:

\begin{table}[H]
\centering\small
\caption{Submodule responsibilities (v0.2.94)}
\resizebox{\textwidth}{!}{
\begin{tabular}{@{}llr@{}}
\toprule
\textbf{Module} & \textbf{Responsibility} & \textbf{LOC} \\
\midrule
\texttt{constants}       & AASM thresholds, band limits, scoring profiles & 76 \\
\texttt{utils}           & Sleep mask construction, channel detection & 134 \\
\texttt{signal}          & Linearisation, baseline, MMSD validation & 327 \\
\texttt{breath}          & Breath segmentation, flattening index & 254 \\
\texttt{classify}        & 7-rule apnea type decision tree & 310 \\
\texttt{spo2}            & SpO$_2$ coupling, ODI, hypoxic burden & 545 \\
\texttt{ecg\_effort}     & ECG-derived effort (TECG + spectral classifier) & 180 \\
\texttt{plm}             & Periodic Limb Movement detection (AASM + WASM) & 271 \\
\texttt{ancillary}       & Heart rate, snoring, position, Cheyne-Stokes & 277 \\
\texttt{signal\_quality} & Per-channel quality grading & 160 \\
\texttt{postprocess}     & CSR reclassification, mixed decomposition, CII & 310 \\
\texttt{respiratory}     & Apnea/hypopnea detection + 12 corrections & 1061 \\
\texttt{pipeline}        & Master orchestrator (11 steps) & 802 \\
\bottomrule
\end{tabular}}
\end{table}

The public API exports 42 symbols from
\texttt{psgscoring/\_\_init\_\_.py}.

% ============================================================
\section{Signal Processing Chain}
% ============================================================

\subsection{Step 1: Nasal Pressure Linearisation}
The nasal pressure transducer measures pressure, not flow. The
Bernoulli relation gives
$x_{\text{lin}}(t) = \text{sign}(x(t)) \cdot \sqrt{|x(t)|}$.

\begin{lstlisting}
def linearise_nasal_pressure(signal):
    return np.sign(signal) * np.sqrt(np.abs(signal))
\end{lstlisting}

\subsection{Step 2: Bandpass Filtering}
A 4th-order Butterworth bandpass (0.05--3~Hz) removes DC drift and
high-frequency noise. The normal respiratory rate (12--20/min =
0.2--0.33~Hz) falls well within this passband.

\subsection{Step 3: Hilbert Envelope}
The instantaneous amplitude: $A(t) = |x(t) + i\mathcal{H}\{x(t)\}|$.
This provides a smooth, non-negative representation of breath amplitude.
The envelope smooths over individual breath cycles---useful for apneas
but can miss subtle hypopneas (motivating peak-based detection in
Section~\ref{sec:undercounting}).

\subsection{Step 4: Dynamic Baseline}
A 5-minute sliding window, 95th percentile, advancing in 10-second
steps with linear interpolation:
$B(t) = P_{95}\{A(\tau) : \tau \in [t-150\text{s},\, t+150\text{s}]\}$.

\subsection{Step 5: Event Detection}
An event is detected when normalised amplitude drops below threshold
for $\geq$10~s:
$A(t)/B(t) \leq 1 - \theta$ where $\theta = 0.90$ for apneas,
$\theta = 0.30$ for hypopneas.

\subsection{Step 6: MMSD Validation}
The Mean Magnitude of Second Derivative
$\overline{|x''(t)|}$ distinguishes true apneas (low but non-zero
MMSD from cardiac oscillation) from artefacts (MMSD $\approx$ 0).

% ============================================================
\section{Apnea Type Classification}
% ============================================================

A 7-rule decision tree classifies each event as obstructive, central,
or mixed based on thoraco-abdominal effort signals.

\subsection{Feature Extraction}
For each event: \texttt{effort\_ratio} (RMS during event / baseline),
\texttt{raw\_var\_ratio}, \texttt{paradox\_corr} (Pearson of thorax vs.
abdomen), \texttt{phase\_angle\_deg} (Hilbert phase difference),
\texttt{quarter\_efforts}, \texttt{flattening\_index}.

\subsection{Decision Tree}

\begin{table}[H]
\centering\small
\caption{Classification rules (first match wins)}
\begin{tabular}{@{}clll@{}}
\toprule
\textbf{Rule} & \textbf{Condition} & \textbf{Type} & \textbf{Conf.} \\
\midrule
0 & Phase angle $\geq 45^{\circ}$ & Obstructive & 0.75--0.97 \\
1 & Paradoxical + high variability & Obstructive & 0.85--0.95 \\
2 & High variability, low envelope & Obstructive & 0.70--0.85 \\
3 & First half low, second half high & Mixed & 0.70--0.95 \\
4 & Clear effort present & Obstructive & 0.65--0.80 \\
5/5a/5b & No effort (+ ECG confirmation) & Central & 0.45--0.90 \\
6 & Default (low effort $\to$ central; else obstr.) & varies & $<$0.60 \\
\bottomrule
\end{tabular}
\end{table}

\textbf{The cardiac pulsation problem:} during a true central apnea,
the heart continues beating. These mechanical oscillations appear on the
RIP bands as low-amplitude signals. Rules~5/5a/5b use relaxed
thresholds and ECG-derived effort to handle this.

% ============================================================
\section{SpO\texorpdfstring{$_2$}{2} Coupling}
% ============================================================

AASM Rule~1A requires $\geq$3\% desaturation. The algorithm computes
the SpO$_2$ baseline (90th percentile during sleep), searches for the
nadir in a post-event window (30--45~s, configurable), and checks
desaturation $\geq$3\%.

\textbf{Circulatory delay:} finger oximetry has 20--40~s delay. The
standard profile uses 45~s; strict uses 30~s.

% ============================================================
\section{Over-Counting Corrections}
\label{sec:overcounting}
% ============================================================

Six mechanisms that inflate AHI:

\textbf{Fix~1: Post-Apnea Baseline Inflation.}
Recovery hyperpnea enters the baseline window, inflating it. Fix:
exclude a 30-second recovery mask after each apnea.

\textbf{Fix~2: SpO$_2$ Cross-Contamination.}
At AHI~$>$60/h, the nadir of event~$N$ overlaps event~$N+1$'s window.
Fix: flag overlapping events (informational since v0.8.14).

\textbf{Fix~3: Cheyne-Stokes Flagging.}
Each CSR trough crosses the 30\% threshold. Fix: detect periodicity
via autocorrelation; flag matching events as \texttt{csr\_flagged}.

\textbf{Fix~4: Confidence Stratification.}
Poor RIP quality defaults to obstructive. Fix: report events at four
confidence tiers ($\geq$0.85, $\geq$0.60, $\geq$0.40, all).

\textbf{Fix~5: Artefact-Flank Exclusion.}
Signal recovery ramp after a flat-line gap. Fix: append 15-second
exclusion window after each gap $\geq$10~s.

\textbf{Fix~6: Local Baseline Validation.}
Rolling baseline inflated by recovery hyperpnea. Fix: compare event
amplitude with 30-second pre-event window; reject if local reduction
$<$20\% (raised to 30\% when surrounding breathing is stable,
CV~$<$0.30).

% ============================================================
\section{Under-Counting Corrections}
\label{sec:undercounting}
% ============================================================

Six mechanisms that suppress detection:

\textbf{Fix~7: Peak-Based Breath Detection.}
AASM requires \emph{peak} excursion $\geq$30\%. The Hilbert envelope
smooths over individual breaths. Fix: detect breaths via zero-crossings,
compute per-breath amplitude, combine with envelope mask via logical OR.
Requires $\geq$3 consecutive reduced breaths (anti-fragmentation).

\textbf{Fix~8: Extended SpO$_2$ Window.}
30~s misses late nadirs. Fix: extend to 45~s (configurable).

\textbf{Fix~9: Flow Smoothing.}
\emph{Removed.} Ablation analysis on PSG-IPA showed that 3-second
smoothing was the dominant source of over-counting (+54 false hypopneas
on a single mild recording). The consecutive-breath requirement
(Fix~7) now replaces smoothing as the anti-fragmentation mechanism.

\textbf{Fix~10: Position Signal Auto-Mapping.}
Some EDF systems store position as raw ADC values (0--255). Fix:
auto-detect encoding; apply percentile-based quantisation.

\textbf{Fix~11: Configurable Scoring Profiles.}
Three profiles balance sensitivity vs. specificity:

\begin{table}[H]
\centering\small
\caption{Scoring profiles}
\begin{tabular}{@{}lccccc@{}}
\toprule
\textbf{Profile} & \textbf{Thresh.} & \textbf{SpO$_2$ win.} &
\textbf{Peak det.} & \textbf{Stab.\ filter} & \textbf{Use case} \\
\midrule
Strict    & $\geq$30\% & 30~s & Envelope only & CV$<$0.45 & Research \\
Standard  & $\geq$30\% & 45~s & Peak+envelope & CV$<$0.45 & Clinical \\
Sensitive & $\geq$25\% & 45~s & Peak+envelope & CV$<$0.45 & UARS \\
\bottomrule
\end{tabular}
\end{table}

\textbf{Fix~12: Flattening-Based RERA Detection.}
The flattening index $FI$ measures the fraction of inspiratory time
exceeding 80\% of peak amplitude. Normal $\approx$0.10; flow-limited
$>$0.30. Sequences of $\geq$3 breaths with FI~$>$0.30 spanning
$\geq$10~s, terminated by an arousal, are scored as flattening-RERAs.
$\text{RDI} = \text{AHI} + (\text{FRI-RERA} + \text{flat-RERA})/\text{TST}$.

% ============================================================
\section{Breath-Amplitude Stability Filter}
% ============================================================

At low event density, normal breath-to-breath variability generates
apparent $\geq$30\% reductions. For each hypopnea candidate, the system
examines breath amplitudes in a $\pm$120-second window. If the
coefficient of variation CV~$<$0.45 with $\geq$10 breaths available,
the candidate is rejected as normal variability.

Calibrated on PSG-IPA: rejected 56 false-positive hypopneas in a
normal recording (SN4: $-$68\%) while preserving accuracy in severe OSA
(SN3: $-$3\%).

% ============================================================
\section{AHI Confidence Interval}
% ============================================================

Every analysis runs all three profiles and reports the AHI as an
interval with a robustness grade:
\textbf{A}~(all profiles same severity),
\textbf{B}~(2/3 concordant),
\textbf{C}~(all discordant).

\begin{lstlisting}
results = run_pneumo_analysis(raw, hypno, channel_map)
iv = results["ahi_interval"]
print(f"AHI: {iv['standard']['ahi']}")
print(f"Interval: [{iv['strict']['ahi']}"
      f"-{iv['sensitive']['ahi']}]")
print(f"Grade: {iv['robustness_grade']}")
\end{lstlisting}

\textbf{Clinical example (SN2):} 12~technologists split 8/4 between
Normal and Mild. Algorithm: interval [4.9--8.9], Grade~B. The interval
immediately communicates that the AHI hovers near the 5/h threshold.

% ============================================================
\section{ECG-Derived Effort Classification}
% ============================================================

\textbf{Spectral classifier:} PSD of the RIP signal during each event.
If $>$75\% of power in the cardiac band (0.8--2.5~Hz) and $<$20\% in
respiratory (0.1--0.5~Hz) $\to$ flag as probable central.

\textbf{TECG (Berry et~al., 2019):} high-pass filter ECG at 30~Hz,
blank QRS ($\pm$50~ms around R-peaks). The residual reveals inspiratory
EMG bursts. If no bursts detected $\to$ no effort $\to$ central
\cite{berry2019}.

% ============================================================
\section{Hypoxic Burden}
\label{sec:hb}
% ============================================================

\subsection{Clinical Motivation}

The AHI counts events but ignores their physiological impact.
Two patients with the same AHI can have radically different
cardiovascular risk depending on the \emph{depth} and \emph{duration}
of their desaturations. The hypoxic burden captures this by integrating
the total oxygen deficit \cite{azarbarzin2019}.

\subsection{Algorithm}

For each respiratory event~$i$:

\begin{enumerate}[nosep]
  \item Compute pre-event SpO$_2$ baseline: 90th percentile in a
    120-second window, with global 95th percentile fallback.
  \item Integrate desaturation area:
    $\text{area}_i = \int_{t_{\text{onset}}}^{t_{\text{recovery}}}
    \max(0,\; B_{\text{SpO}_2} - \text{SpO}_2(t))\, dt$
  \item Normalise:
    $\text{HB} = \sum_i \text{area}_i \;/\; (60 \times \text{TST}_{\text{hours}})$
    \quad [\%$\cdot$min/h]
\end{enumerate}

\begin{table}[H]
\centering\small
\caption{Hypoxic burden risk thresholds (Azarbarzin et~al.\ 2019)}
\begin{tabular}{@{}lcl@{}}
\toprule
\textbf{Risk} & \textbf{HB (\%$\cdot$min/h)} & \textbf{Clinical} \\
\midrule
Low      & $<$20  & Low CV risk \\
Moderate & 20--73 & Moderate risk \\
High     & $>$73  & HR 1.47 for incident CVD \\
\bottomrule
\end{tabular}
\end{table}

\begin{lstlisting}
from psgscoring import compute_hypoxic_burden
result = compute_hypoxic_burden(spo2_data, sf_spo2=1.0,
                                resp_events=events, hypno=hypno)
print(f"HB: {result['hypoxic_burden']} {result['unit']}")
\end{lstlisting}

\subsection{Baseline Methods}

Two methods for computing the per-event SpO$_2$ baseline are available
via the \texttt{baseline\_method} parameter:

\textbf{Percentile method} (\texttt{"percentile"}, default).
Uses the 90th percentile of the 120-second pre-event window as
baseline, with the global 95th percentile of all sleep SpO$_2$ as
fallback (whichever is higher). Simple, robust, and validated by
He et~al.\ (IEEE EMBC 2024) as comparable to the original ensemble
method for cardiovascular mortality prediction.

\textbf{Ensemble method} (\texttt{"ensemble"}, Azarbarzin original).
All SpO$_2$ segments are synchronised at event termination (``time
zero'') and averaged to produce a subject-specific ensemble curve.
The nadir of this curve is detected, and the two flanking peaks define
a search window.  The baseline for each individual event is the SpO$_2$
at the left peak; the area is integrated within the search window.
Falls back to percentile if $<$3 events are available.

\begin{lstlisting}
# Percentile (default, backward compatible)
hb1 = compute_hypoxic_burden(spo2, sf, events, hypno)

# Ensemble (Azarbarzin 2019 original)
hb2 = compute_hypoxic_burden(spo2, sf, events, hypno,
                              baseline_method="ensemble")
print(f"Ensemble window: {hb2['ensemble_window_s']} s")
\end{lstlisting}

On synthetic data (30 events, 8-hour recording), the two methods
differ by approximately 8\%.  The ensemble method tends to yield
slightly higher values because the search window can extend beyond
the recovery point used by the percentile method.  For the MESA
validation, both methods will be compared on $\sim$2,000 recordings.

\subsection{Validation: The Pevernagie Phenotype}

On PSG-IPA, SN2 (AHI~8.9, Mild) has HB~71.7---approaching the
high-risk threshold. SN1 (AHI~7.6, also Mild) has HB~12.5. Same
AHI category, 6$\times$ difference in oxygen deficit. On iSLEEPS
($n=96$), 9~patients with AHI~$<$15 had HB~$>$40, including one with
AHI~9.0 and HB~114.

% ============================================================
\section{Post-Processing: Central Apnea Refinement}
\label{sec:postprocess}
% ============================================================

After the main detection pipeline, three refinements are applied
automatically (pipeline Step~11).

\subsection{CSR-Aware Central Reclassification}

During Cheyne-Stokes respiration, cardiac pulsation on the effort
bands mimics respiratory effort, causing central apneas to be
misclassified as obstructive. Events with \texttt{csr\_flagged=True}
that are typed \texttt{"obstructive"} or \texttt{"mixed"} are
reclassified as \texttt{"central"}.

Validation: PSG-IPA: 23~events reclassified. iSLEEPS ($n=96$):
365~events in 36/96 studies (37\%).

\subsection{Mixed Apnea Decomposition}

Each mixed apnea is decomposed into central (absent effort) and
obstructive (resumed effort) portions. If the central portion
$\geq$10~s, the event is reclassified as central.

\begin{lstlisting}
from psgscoring import decompose_mixed_apneas
events = decompose_mixed_apneas(events, thorax, abdomen, sf)
# Each mixed event gains: central_duration_s,
#   obstructive_duration_s, central_ratio
\end{lstlisting}

\subsection{Central Instability Index (CII)}

Quantifies profile-dependent O/C uncertainty by comparing AHI across
profiles. CII~$<$0.15 = stable; 0.15--0.40 = moderate;
$>$0.40 = high (manual review recommended).

% ============================================================
\section{Arousal Detection}
% ============================================================

\textbf{AASM criteria:} abrupt EEG frequency shift ($\geq$3~s, preceded
by $\geq$10~s stable sleep; during REM: concurrent EMG increase).

Phase~1 (EEG-based): sliding-window alpha+beta power exceeding
mean~$+$~2$\times$SD. Phase~2 (filtering): K-complex exclusion (raises
minimum duration from 3~s to 5~s), CVR confidence boost
($\geq$10~bpm swing within 15~s), respiratory coupling (within 3~s of
event termination $\to$ respiratory arousal), Rule~1B reinstatement
(flow reduction + arousal without desaturation $\to$ hypopnea).

% ============================================================
\section{Ancillary Analyses}
% ============================================================

\textbf{Snoring:} bandpass 30--300~Hz, 1-second RMS windows, threshold
at P60 of RMS during sleep.

\textbf{Heart rate:} R-peak detection via bandpass 5--30~Hz +
\texttt{find\_peaks()}. Metrics: mean/min/max HR, bradycardia/tachycardia
episodes.

\textbf{Body position:} auto-mapping of ADC/integer/string encodings.
Positional AHI per position. Supine-dominant OSA: supine~AHI~$\geq
2\times$~non-supine. Position-induced baseline reset at transitions.

\textbf{Cheyne-Stokes:} autocorrelation of flow envelope over
5-minute windows; secondary peak at 40--100~s confirms CSR periodicity.

\textbf{PLM:} AASM~2.6 + WASM criteria (0.5--10~s duration, $\geq$8~$\mu$V,
5--90~s interval, $\geq$4 consecutive). Respiratory-associated movements
excluded from PLMI.

\textbf{Sleep cycles:} Feinberg \& Floyd criteria (NREM $\geq$15~min,
REM consolidation with 2-min gap tolerance). Stage-specific AHI (REM
vs.\ NREM) identifies REM-dominant phenotype.

% ============================================================
\section{Signal Quality Assessment}
% ============================================================

Every channel is evaluated on four dimensions: flat-line percentage
(rolling 2-second SD), clipping ($>$10 consecutive samples at ADC rail),
line-noise dominance (49--51 and 59--61~Hz), and disconnect events.
Each channel receives a grade (good/acceptable/poor/unusable).

Montage plausibility: EEG$\leftrightarrow$EOG correlation $>$0.95
suggests shared reference; thorax$\leftrightarrow$abdomen $>$0.98
suggests duplication.

% ============================================================
\section{The Complete Pipeline}
% ============================================================

\texttt{run\_pneumo\_analysis()} orchestrates 11 steps:

\begin{table}[H]
\centering\small
\caption{Pipeline steps}
\begin{tabular}{@{}cl@{}}
\toprule
\textbf{Step} & \textbf{Description} \\
\midrule
1  & Signal quality assessment \\
2  & Linearisation + filtering + envelope \\
3  & Dynamic baseline computation \\
4  & Apnea detection (thermistor, $\geq$90\%) \\
5  & Hypopnea detection (nasal pressure, $\geq$30\%) \\
6  & SpO$_2$ coupling + type classification \\
7  & 6 over-counting corrections \\
8  & 6 under-counting corrections + stability filter \\
9  & AHI confidence interval (3 profiles) \\
10 & Hypoxic burden (Azarbarzin 2019) \\
11 & Post-processing (CSR recl.\ + mixed decomp.\ + CII) \\
\bottomrule
\end{tabular}
\end{table}

% ============================================================
\section{External Validation}
% ============================================================

\subsection{PSG-IPA (5 recordings, 60 scorer sessions)}

\begin{table}[H]
\centering\small
\caption{PSG-IPA: algorithm vs.\ 12 scorers per recording}
\resizebox{\textwidth}{!}{
\begin{tabular}{@{}lrrccccrl@{}}
\toprule
& \textbf{$n$} & \textbf{AHI\textsubscript{med}} &
\textbf{AHI\textsubscript{algo}} & \textbf{$|\Delta|$} &
\textbf{O} & \textbf{C} & \textbf{HB} & \textbf{CSR recl.} \\
\midrule
SN1 & 11 &  6.0 &  7.6 & 1.6 &  9 &  7 & 12.5 & 7 \\
SN2 & 12 &  3.7 &  8.9 & 5.2 &  3 &  0 & \textbf{71.7} & 0 \\
SN3 & 12 & 54.0 & 53.5 & 0.5 & 252 & 30 & \textbf{241.8} & 16 \\
SN4 & 12 &  3.5 &  4.2 & 0.7 &  0 &  0 & 22.1 & 0 \\
SN5 & 12 &  9.9 & 10.7 & 0.8 & 17 &  1 & 37.0 & 0 \\
\midrule
\multicolumn{4}{@{}l}{Mean $|\Delta|$: 1.8/h} &  && & & 23 total \\
\bottomrule
\end{tabular}}
\end{table}

\noindent Event-level matching on SN3 (severe): F1$=0.890$, mean onset
difference 2.3~s. For 3/5~recordings, algorithm deviation was smaller
than inter-scorer variability.

\subsection{iSLEEPS (96 ischemic stroke patients)}

\begin{table}[H]
\centering\small
\caption{iSLEEPS by severity ($n=96$, 1 outlier excluded)}
\begin{tabular}{@{}lccccc@{}}
\toprule
& \textbf{Overall} & \textbf{Normal} & \textbf{Mild} &
\textbf{Mod.} & \textbf{Severe} \\
& ($n=96$) & ($n=12$) & ($n=20$) & ($n=26$) & ($n=38$) \\
\midrule
Bias   & $-$14.3/h & $+$1.3 & $-$1.3 & $-$12.4 & $-$27.5 \\
MAE    & 15.6/h & 3.4 & 3.9 & 13.0 & 27.5 \\
$r$    & 0.653 &&&&\\
Conc.  & 30\% & 83\% & 45\% & 23\% & 11\% \\
\bottomrule
\end{tabular}
\end{table}

\noindent 365 CSR reclassifications in 36/96 studies. 15~patients with
HB~$>$73 (high CV risk). 9~patients with AHI~$<$15 but HB~$>$40.

% ============================================================
\section{Data Structures}
% ============================================================

\subsection{Event Dictionary}
\begin{lstlisting}
{
    "type": "obstructive",  # or central/mixed/hypopnea
    "onset_s": 1234.5,
    "duration_s": 18.0,
    "stage": "N2",
    "desaturation_pct": 4.2,
    "min_spo2": 89.3,
    "confidence": 0.82,
    "rule_index": 1,
    "flow_reduction": 0.45,
    "classify_detail": { ... },
    "csr_flagged": False,
    "spo2_cross_contaminated": False,
    # Post-processing (v0.2.94):
    "csr_reclassified": True,
    "original_type": "obstructive",
    "central_duration_s": 12.1,
    "obstructive_duration_s": 7.9,
    "central_ratio": 0.60,
    "mixed_decomposed": True,
}
\end{lstlisting}

\subsection{Hypoxic Burden Output}
\begin{lstlisting}
results["hypoxic_burden"] = {
    "hypoxic_burden": 71.7,       # %*min/h
    "n_events_with_burden": 40,
    "mean_event_burden": 645.3,   # %*s per event
    "unit": "%*min/h",
}
\end{lstlisting}

\subsection{Post-Processing Output}
\begin{lstlisting}
results["postprocess"] = {
    "n_csr_reclassified": 7,
    "n_mixed_decomposed": 3,
    "n_mixed_to_central": 2,
    "central_instability": {
        "central_instability_index": 0.43,
        "interpretation": "high",
        "ahi_range": 3.3,
    },
    "cai_change": +9,
}
\end{lstlisting}

% ============================================================
\section{Practical Tutorial}
% ============================================================

\begin{lstlisting}
import mne
from psgscoring import run_pneumo_analysis

raw = mne.io.read_raw_edf("patient.edf", preload=True)
hypno = ["W","N1","N2","N3","R", ...]  # per 30-s epoch

results = run_pneumo_analysis(raw, hypno, channel_map={
    "flow_pressure": "Pressure",
    "flow_thermistor": "Flow",
    "thorax": "RIP Thora", "abdomen": "RIP Abdom",
    "spo2": "SpO2", "ecg": "ECG II",
}, scoring_profile="standard")

# AHI with confidence interval
iv = results["ahi_interval"]
print(f"AHI: {iv['standard']['ahi']}")
print(f"[{iv['strict']['ahi']}"
      f"-{iv['sensitive']['ahi']}] "
      f"Grade {iv['robustness_grade']}")

# Hypoxic burden
hb = results["hypoxic_burden"]
print(f"HB: {hb['hypoxic_burden']} {hb['unit']}")

# Post-processing
pp = results["postprocess"]
print(f"CSR reclassified: {pp['n_csr_reclassified']}")
\end{lstlisting}

% ============================================================
\section{Design Decisions}
% ============================================================

\textbf{Why rules, not deep learning?} Transparency (every event
traceable to AASM rule), data efficiency (zero training data),
auditability (CE/FDA), correction specificity (each mechanism
independently addressable). The planned hybrid uses LightGBM
\emph{on top of} the rule-based engine.

\textbf{Performance:} 3--5~min for an 8-hour PSG at 256~Hz.
All 12 corrections add $<$1~s overhead.

% ============================================================
\section{Common Pitfalls}
% ============================================================

\textbf{Variable scoping:} pass computed variables explicitly to inner
functions---Python closures capture by reference.
\textbf{Docker caching:} \texttt{docker compose restart} does not
rebuild; use \texttt{build --no-cache}.
\textbf{Shell escaping:} use Python parameterised queries for
passwords (bcrypt \$ interpretation).
\textbf{EDF:} pyedflib for EDF+ writing; edfio silently drops
annotations.
\textbf{Translation shadowing:} never \texttt{t = date.today()} when
\texttt{t()} is imported.

% ============================================================
\section{Exercises}
% ============================================================

\begin{enumerate}
  \item Generate a synthetic nasal pressure signal (12 breaths/min).
    Apply square-root linearisation. Compare the envelope before and
    after.
  \item Create a 10-minute signal with a 30-second apnea and 150\%
    recovery breathing. Plot the dynamic baseline with and without
    the 30-second recovery mask (Fix~1).
  \item Create test inputs for \texttt{classify\_apnea\_type()} that
    trigger each of the 7 rules.
  \item Generate a SpO$_2$ signal with a 5\% desaturation at $t=40$~s.
    Test \texttt{get\_desaturation()} with 30~s and 45~s windows.
  \item Download a public PSG from PhysioNet. Run all three scoring
    profiles and compare AHI.
  \item Generate a SpO$_2$ signal with three desaturations (3\%, 8\%,
    15\%). Call \texttt{compute\_hypoxic\_burden()} and verify the deep
    event dominates.
  \item Create 10 events with \texttt{csr\_flagged=True} and type
    \texttt{"obstructive"}. Call \texttt{reclassify\_csr\_events()} and
    verify reclassification.
  \item Generate a 20-second effort signal: 12~s silence then 8~s
    sinusoidal. Call \texttt{decompose\_mixed\_apneas()} and verify
    central reclassification.
  \item Using the iSLEEPS CSV, identify patients where AHI~$<$15 but
    HB~$>$40. How does hypoxic burden change the clinical picture?
\end{enumerate}

% ============================================================
\section*{Glossary}

\begin{tabular}{@{}lp{10cm}@{}}
AHI & Apnea-Hypopnea Index \\
OAHI & Obstructive AHI \\
CAI & Central Apnea Index \\
RDI & Respiratory Disturbance Index (AHI + RERA index) \\
RERA & Respiratory Effort-Related Arousal \\
ODI & Oxygen Desaturation Index \\
HB & Hypoxic Burden (\%$\cdot$min/h) \\
CII & Central Instability Index \\
PLMI & Periodic Limb Movement Index \\
TST & Total Sleep Time \\
RIP & Respiratory Inductance Plethysmography \\
TECG & Transformed ECG (high-pass + QRS-blanked) \\
CSR & Cheyne-Stokes Respiration \\
MMSD & Mean Magnitude of Second Derivative \\
EDF & European Data Format \\
AASM & American Academy of Sleep Medicine \\
\end{tabular}

% ============================================================
\begin{thebibliography}{29}

\bibitem{berry2020}
Berry RB, et~al. \emph{The AASM Manual for the Scoring of Sleep and
Associated Events. Version~2.6.} AASM; 2020.

\bibitem{vallat2021}
Vallat R, Walker MP. An open-source, high-performance tool for
automated sleep staging. \emph{eLife}. 2021;10:e70092.

\bibitem{pevernagie2020}
Pevernagie DA, et~al. On the rise and fall of the apnea--hypopnea
index: a historical review and critical appraisal. \emph{J Sleep Res}.
2020;29(4):e13066.

\bibitem{pevernagie2025}
Pevernagie DA, et~al. Sleep Medicine---What's in a Name?
\emph{J Sleep Res}. 2025;34(5):e70150.

\bibitem{azarbarzin2019}
Azarbarzin A, et~al. The hypoxic burden of sleep apnoea is an
independent predictor of incident cardiovascular outcomes.
\emph{Am J Respir Crit Care Med}. 2019;200(2):211--219.

\bibitem{berry2019}
Berry RB, et~al. Use of a transformed ECG signal to detect respiratory
effort during apnea. \emph{J Clin Sleep Med}. 2019;15(11):1653--1660.

\bibitem{berry2016}
Berry RB, et~al. Use of chest wall electromyography to detect
respiratory effort. \emph{J Clin Sleep Med}. 2016;12(9):1239--1244.

\bibitem{rosenberg2013}
Rosenberg RS, Van Hout S. The AASM inter-scorer reliability program.
\emph{J Clin Sleep Med}. 2013;9(1):81--87.

\bibitem{aasm_isr2014}
Rosenberg RS, et~al. The AASM inter-scorer reliability program:
respiratory events. \emph{J Clin Sleep Med}. 2014;10(4):447--454.

\bibitem{ruehland2009}
Ruehland WR, et~al. The new AASM criteria for scoring hypopneas.
\emph{Sleep}. 2009;32(2):150--157.

\bibitem{malhotra2013}
Malhotra A, et~al. Performance of an automated polysomnography
scoring system. \emph{Sleep}. 2013;36(4):573--582.

\bibitem{parekh2023}
Parekh A, et~al. Ventilatory burden as a measure of OSA severity.
\emph{Am J Respir Crit Care Med}. 2023;208(11):1216--1226.

\bibitem{thurnheer2001}
Thurnheer R, et~al. Accuracy of nasal cannula pressure recordings.
\emph{Am J Respir Crit Care Med}. 2001;164(10):1914--1919.

\bibitem{hosselet1998}
Hosselet J, et~al. Detection of flow limitation with a nasal
cannula/pressure transducer system. \emph{Am J Respir Crit Care Med}.
1998;157(5):1461--1467.

\bibitem{lee2008}
Lee H, et~al. Detection of apneic events from single-channel nasal
airflow. \emph{Physiol Meas}. 2008;29:N37--N45.

\bibitem{perslev2021}
Perslev M, et~al. U-Sleep: resilient high-frequency sleep staging.
\emph{npj Digital Medicine}. 2021;4:72.

\bibitem{lightgbm}
Ke G, et~al. LightGBM: A highly efficient gradient boosting decision
tree. \emph{NeurIPS}. 2017;30:3146--3154.

\bibitem{gramfort2014}
Gramfort A, et~al. MNE software for processing MEG and EEG data.
\emph{NeuroImage}. 2014;86:446--460.

\bibitem{feinberg1979}
Feinberg I, Floyd TC. Systematic trends across the night in human
sleep cycles. \emph{Psychophysiology}. 1979;16(3):283--291.

\bibitem{zucconi2006}
Zucconi M, et~al. WASM standards for recording and scoring PLM.
\emph{Sleep Med}. 2006;7(2):175--183.

\bibitem{peppard2013}
Peppard PE, et~al. Increased prevalence of sleep-disordered breathing.
\emph{Am J Epidemiol}. 2013;177(9):1006--1014.

\end{thebibliography}

\vspace{1cm}
\noindent\texttt{pip install psgscoring}\\
\url{https://github.com/bartromb/psgscoring}\\
\url{https://slaapkliniek.be}

\end{document}
